\section{Conclusion}
\label{sec: conclusion}


Major advances in deep learning have resulted from powerful architectural innovations enabling previously-handcrafted features to be learned from data,
from CNNs learning visual features~\citep{AlexNet} to Transformers discovering linguistic patterns~\citep{Transformer}.
\textbf{H-Nets} similarly unlock the ability to remove another layer of pre-processing, such as tokenizers, and instead learn them end-to-end.
This ability results from a set of new techniques we introduce that work together to form a \textbf{dynamic chunking} mechanism,
which is able to learn content- and context- dependent discrete segmentation strategies through standard gradient-based optimization.
A single-stage byte-level H-Net already exceeds the performance of standard tokenized language models,
and recursive H-Nets with multiple stages of dynamic chunking further improve its scaling.
H-Nets substantially remedy issues with tokenizers, display very strong performance on diverse languages and language-like modalities,
and more broadly may serve as the backbone of general foundation models that do \emph{more learning} with \emph{less processing}.

